\begin{esercizio}{Acqua nel contenitore graduato}
In laboratorio, una studentessa riempie un contenitore graduato con $12\unita{dl}$ d'acqua.

\begin{enumerate}[label=\alph*)]
  \item Qual è la capacità dell'acqua contenuta nel recipiente, espressa in litri?

  \spazioRisposta[1.5cm]

  \item Qual è il volume dell'acqua contenuta nel recipiente, espresso in $\unita{dm^3}$?

  \spazioRisposta[1.5cm]

  \item Qual è il volume dell'acqua contenuta nel recipiente, espresso in $\unita{cm^3}$?

  \spazioRisposta[1.8cm]
\end{enumerate}
\end{esercizio}

\begin{esercizio}{Equivalenze}
Completa le seguenti equivalenze.

\begin{enumerate}[label=\alph*)]
  \item $3\unita{l} = \rule{2cm}{0.4pt}\unita{dl}$
  \item $45\unita{dl} = \rule{2cm}{0.4pt}\unita{l}$
  \item $2.5\unita{l} = \rule{2cm}{0.4pt}\unita{ml}$
  \item $750\unita{ml} = \rule{2cm}{0.4pt}\unita{l}$
  \item $4\unita{dm^3} = \rule{2cm}{0.4pt}\unita{l}$
  \item $320\unita{cm^3} = \rule{2cm}{0.4pt}\unita{ml}$
\end{enumerate}
\end{esercizio}

\begin{esercizio}{Problema}
Una bottiglia contiene $1.5\unita{l}$ d'acqua. Matteo beve $350\unita{ml}$ d'acqua.

\begin{enumerate}[label=\alph*)]
  \item Quanti millilitri d'acqua conteneva inizialmente la bottiglia?

  \spazioRisposta[1.3cm]

  \item Quanti millilitri d'acqua rimangono nella bottiglia?

  \spazioRisposta[1.3cm]

  \item Esprimi la quantità rimasta in litri.

  \spazioRisposta[1.5cm]
\end{enumerate}
\end{esercizio}
